\section{Related Work}
\label{sec:related_work}

This section reviews five interconnected areas that motivate the
design of the MVC Test Orchestrator. It first examines traditional
approaches to architectural enforcement and static analysis,
highlighting their structural limitations. It then surveys the
evolution of LLM-based code generation and its implications for
architectural integrity. Next, it reviews requirement-to-architecture
traceability frameworks and their limitations in supporting automated,
end-to-end compliance. It also assesses automated program repair and
code remediation, and finally discusses Human-AI collaboration in
software engineering. Together, these strands reveal the specific gap
the proposed system is designed to address.

% -----------------------------------------------------------------
\subsection{Traditional Architectural Enforcement and Static Analysis}
\label{sec:related_static}
% -----------------------------------------------------------------

Software architecture serves as a foundational building block for
satisfying a system's quality attributes and enforcing its structural
constraints. As software systems and their component modules evolve,
the separation between the intended architectural design and the
implemented codebase becomes an ongoing challenge. Defined in the
literature as Architecture Erosion \citep{Li2022}, this
deviation risks the integrity of the system's modular structure.
Beyond structural degradation, it increases system vulnerabilities
and limits extensibility, complicating the integration of new
requirements \citep{Li2022}. Although various monitoring and
remediation strategies have been proposed to mitigate this decay
\citep{deSilva2012}, the interplay of technical and non-technical
factors keeps architecture erosion among the major unsolved problems in software
engineering \citep{Li2022}.

Historically, approaches to preserving architectural integrity have
relied on modularization, information hiding \citep{Parnas1972}, and
optimized coupling and cohesion \citep{Yourdon1979}.
Scenario-based evaluation frameworks, including SAAM \citep{Kazman1994}
and ATAM \citep{Kazman2000}, were developed to identify architectural
risks before implementation and to verify alignment between quality
requirements and the proposed design. However, most traditional methods
fail to establish a consistent link between quality attributes and
architectural artifacts, and often depend on manual or informal
processes \citep{Dobrica2002}. In practice, failure to strictly capture
high-level architectural artifacts---documents that represent the
macro-level system design---often prevents these abstractions from
guiding the implementation and reuse of low-level components
\citep{Lung1997}.

Currently, automated tools used to ensure architectural compliance,
such as SonarQube \citep{SonarQube} and Checkstyle
\citep{Checkstyle}, focus almost exclusively on low-level code
anomalies, including syntactic irregularities and stylistic errors.
Consequently, high-level structural violations affecting the system's
organizational logic---such as invalid dependencies that violate
architectural rules---are detected only to a limited extent \citep{Fontana2017}.
Developers' reluctance to adopt formal architectural compliance tools
has been attributed in the literature to the substantial effort
required to map a system to its architecture \citep{Ali2018}. Existing commercial solutions typically
require substantial manual preprocessing to map the codebase to the
intended architecture---a configuration overhead that \citet{Ali2018} identify as a key adoption barrier, often resulting in
informal, error-prone verification methods.

Operating under a ``Just-In-Time'' paradigm, tools like JITTAC
\citep{Buckley2013} represent some of the most proactive traditional
solutions, as developers are informed about the architectural
implications of their modifications during coding. By employing
real-time Reflexion Modeling \citep{Murphy2001} techniques, they reduce
architectural drift by enforcing consistency between the designed
architecture and the executed code. Nevertheless, even with
sophisticated tools like JITTAC, the target architectural model must
still be manually defined by software architects, and source code
elements must be explicitly mapped to this abstraction
\citep{Buckley2013}.

Recent studies \citep{Fuchs2025, Helmi2025} show that LLM-based
approaches have significant potential to autonomously extract and
isolate architectural components from Software Requirement
Specifications (SRS), thereby addressing manual mapping bottlenecks.
To reduce the risk of architectural drift, however, bidirectional
traceability between documented requirements and code-level structural
constraints should be coupled with autonomous remediation.

Architecture conformance approaches differ from general code-quality analyzers such as SonarQube~\citep{SonarQube} and Checkstyle~\citep{Checkstyle}. Reflexion models~\citep{Murphy2001} and just-in-time conformance techniques~\citep{Buckley2013} compare implemented dependencies with an explicitly defined target architecture, while architecture-test and dependency-rule approaches encode admissible relations as executable constraints~\citep{Pruijt2017}. These families provide stronger architectural baselines than style-oriented analyzers, but typically require developers to specify or maintain the target model. The present prototype explores whether an SRS-derived MVC map can reduce that setup burden; it does not replace or empirically outperform established conformance tools in this study.

The tools and methods reviewed in this section share a common
limitation: they either operate at too low an abstraction level to
capture MVC-layer violations or require substantial manual
configuration, undermining automation. The MVC Test Orchestrator
addresses this gap by combining SRS-grounded requirement analysis with
automated architectural auditing, eliminating the need for manual
model definition while maintaining developer management through a
Human-in-the-Loop design. Unlike JITTAC and similar tools, which
require architects to manually define the target model before analysis
begins, the MVC Test Orchestrator derives the model automatically from
the SRS document. Detection is also scoped to MVC-specific inter-layer
violations rather than general code anomalies, removing the manual
configuration step identified as a key adoption barrier by \citet{Ali2018}.

% -----------------------------------------------------------------
\subsection{Large Language Models in Software Engineering}
\label{sec:related_llm}
% -----------------------------------------------------------------

The integration of Large Language Models (LLMs) into the software
development lifecycle has driven a major shift in automated code
generation and developer productivity. Despite the well-documented
risks to logical consistency and contextual accuracy \citep{Siddiq2024},
the speed of these models has supported their general adoption across
the industry \citep{Hou2024}. Nevertheless, \citet{Slater2025} reports that smaller, open-weight
LLMs exhibit a critical limitation in their ability to naturally
understand and enforce global architectural constraints that control
the macro-level system, whereas larger proprietary models demonstrate
markedly higher conformance.

Current research suggests that while LLMs excel at generating
syntactically correct and functionally feasible code snippets within a
localized context, they may systematically violate layered architectural
rules and isolation principles at the macro level
\citep{Slater2025, Zhu2026}. Particularly within ``vibe coding'' workflows, AI-generated
implementations that conflict with the project's architectural
blueprints are reported to accelerate the accumulation of
Architecture Erosion and technical debt
\citep{Slater2025, Waseem2025, Ehsani2026}. Developers' over-reliance on
AI-generated outputs often allows these structural violations to enter
the codebase without strict architectural inspection, fueling an
uncontrolled accumulation of technical debt \citep{Huang2026}.

Consequently, prompt-level engineering or simple fine-tuning alone may
be insufficient to preserve architectural integrity
\citep{Dente2026, Zhu2026}. This gap underscores the need for an
``architectural orchestrator'' mechanism, one capable of constraining
the generative output rate of LLMs by establishing a structured,
bidirectional bridge between Software Requirement Specifications (SRS)
and code-level structural constraints, while also including autonomous
remediation capabilities.

The MVC Test Orchestrator directly addresses this need by imposing
architectural constraints before code generation begins, grounding each
generation step in SRS-derived requirements rather than relying on the
model to infer architectural rules independently. This structural
approach to constraint enforcement, discussed further in Section~\ref{sec:experiment},
represents a departure from existing LLM coding assistants that lack a
formal architectural representation \citep{Cervantes2026}.

% -----------------------------------------------------------------
\subsection{Requirement-to-Architecture Traceability and RAG}
\label{sec:related_traceability}
% -----------------------------------------------------------------

Traceability, defined as the process of establishing and maintaining
bidirectional relationships between software requirements,
architectural artifacts, and source code, is a critical determinant of
product quality and effective change management \citep{Li2022}. A
recent systematic literature review spanning 2013 to 2024 indicates
that, while most traceability methods rely heavily on Information
Retrieval (IR) techniques, these approaches still exhibit severe
limitations in full automation and seamless integration with modern,
iterative development approaches such as Agile and DevOps
\citep{Koboyatshwene2025}. Furthermore, the scarcity of comprehensive,
end-to-end tool support and the tendency to evaluate these methods on
static, legacy datasets rather than dynamic industrial case studies
continue to slow industrial adoption and limit empirical reliability.

The built-in semantic gap between high-level architectural
documentation and low-level source code constitutes a fundamental
obstacle to establishing robust traceability links. Traditionally,
bridging this gap required the manual construction of software
architectural models---a highly labor-intensive, error-prone, and
costly effort \citep{Clements2002}. However, recent advances in Large
Language Models (LLMs) have enabled the automated extraction of
architectural entities and their relationships directly from both
natural-language documentation and source code. By avoiding the need
for manual modeling, this capability effectively bridges the semantic
gap and enables the generation of high-fidelity traceability links
\citep{Fuchs2025}.

These advances provide a solid foundation for the use of
Retrieval-Augmented Generation (RAG) architectures in software
engineering \citep{Hey2025}. In this context, RAG utilizes autonomously
extracted architectural components and formal requirement
specifications as a semantic ``grounding'' mechanism. This enables
thorough verification of code-level architectural integrity against the
system's dynamic documentation (SRS documents updated
iteratively). Ultimately, this integration harmonizes the generative
capabilities of LLMs with defined architectural constraints,
effectively transforming traceability from a passive, backward-looking
tracking mechanism into an active, real-time remediation process.

The MVC Test Orchestrator builds directly on these advances by
embedding RAG-based semantic retrieval as its primary grounding
mechanism. Rather than treating traceability as a passive
record-keeping activity, the proposed system uses SRS-derived
requirements as a live semantic anchor against which generated code is
continuously verified, closing the automation gap identified in current
IR-based approaches. In this way, automated forward traceability from requirements to code---a
gap that existing methods fail to close---is achieved through the combination
of LLM-based extraction and structured architectural auditing
\citep{dim2026mvc}.

% -----------------------------------------------------------------
\subsection{Automated Program Repair and Code Remediation}
\label{sec:related_apr}
% -----------------------------------------------------------------

While traditional debugging and maintenance methodologies mainly focus
on generating localized patches driven by specific test scenarios, the
advent of generative AI has advanced this area toward the
conceptualization of ``self-healing'' software systems. This
transition forms the core of modern Automated Program Repair (APR)
frameworks, such as RepairAgent \citep{Bouzenia2025} and ChatRepair
\citep{Xia2024}, which are engineered to minimize maintenance overhead
and substantially enhance system reliability. Recent empirical work by
\citet{Khlaisamniang2023} demonstrated that
autonomous remediation mechanisms can effectively address
component-level failures in large-scale systems by leveraging
generative AI for anomaly detection, code generation, and
automated debugging. Such approaches significantly reduce the
need for manual intervention while preserving core system
functionality and operational efficiency.

However, while the automated remediation of localized, syntax-level
defects has become increasingly viable, resolving macro-level
structural debt and architectural ``smells'' \citep{Fontana2017}
remains a highly complex challenge. A recent investigation by
\citet{Pandini2025} highlighted that LLMs exhibit considerable
capability in formulating refactoring strategies to mitigate complex
structural flaws, such as Cyclic Dependencies (CD). Their study
revealed that RAG-enhanced LLM outputs were preferred by 73\% of
participants in qualitative interviews, primarily because they offered
improved solution structure and clearer step-by-step reasoning in the
generated refactoring proposals.

To handle the contextual complexity present in architectural
remediation, the use of multi-agent architectures has been suggested to be
important for preserving system integrity
\citep{Oueslati2025, Xu2025, Batole2025}. The systematic division of
responsibilities---for instance, designating one agent to identify the
optimal dependency to separate, while another executes the
corresponding code refactoring---yields more reliable outcomes when
resolving complex dependency cycles \citep{Pandini2025}. These findings
suggest that coordinated LLM-driven remediation holds promise for
addressing not merely functional anomalies, but also deep structural
deviations that violate a project's foundational architectural
constraints; however, current benchmarks indicate this capability
remains inconsistent, with even the best-performing agents resolving
fewer than half of identified architectural smells and occasionally
introducing new violations \citep{Dinu2026}.

Existing APR frameworks, such as RepairAgent and ChatRepair, operate
without a fixed architectural reference, which means a local fix can
introduce a new structural violation elsewhere in the codebase. The
Fixer Agent in the MVC Test Orchestrator avoids this by restricting
all repairs to the violations flagged by the AST-based static
analyzer; it cannot modify the rule-based project scaffold or the Architectural Map
during patching. The practical effect of this constraint is measured
in Section~\ref{sec:results} through the RSDA and RSR\textsubscript{e2e} metrics.

% -----------------------------------------------------------------
\subsection{Human-AI Collaboration in Software Engineering}
\label{sec:related_hitl}
% -----------------------------------------------------------------

The deployment of AI agents within software engineering workflows is
increasingly shifting away from fully autonomous ``black-box''
approaches and toward Human-in-the-Loop (HITL) hybrid systems. \citet{Sera2025} emphasize that within dynamic,
self-managing architectures, human oversight during the ``Plan''
phase of the classical Monitor, Analyze, Plan, Execute,
Knowledge (MAPE-K) cycle \citep{Kephart2003} is essential for safeguarding
quality attributes and reducing unintended behavior. This
collaborative model positions AI agents as autonomous executors of
operational tasks while reserving strategic governance to human
developers, ensuring the system does not unintentionally deviate from
its intended architectural constraints.

This architectural imperative is supported by current Human-in-the-Loop
LLM-based frameworks. Existing literature demonstrates that permitting
AI agents to execute complex planning and coding tasks without human
validation significantly increases the risk of generating convincing
but erroneous outputs (hallucinations) \citep{Ji2023}.
Ultimately, the effectiveness of human-AI collaboration is dependent
on the transparency and explainability of the AI's recommendations.
Empirical analyses by \citet{Pandini2025} indicate that
developers adopt AI-generated architectural refactoring suggestions
with significantly greater confidence when the proposals are
accompanied by explicit, step-by-step reasoning, thereby directly
enhancing the perceived utility of the intervention.

Building on this collaborative context, the authors initially
introduced the foundational baseline of the MVC Test
Orchestrator~\citep{dim2026mvc}. The preliminary architecture presented at the
International Informatics and Software Engineering Conference (IISEC) constrained LLM outputs using structured, rule-based
scaffolding, establishing the SRS document as the definitive source of
truth. The present study extends this initial framework, as detailed in
Section~\ref{sec:introduction}. Through these advancements, human-AI
interaction is elevated from a passive validation checkpoint into a
proactive, autonomous, and architecturally grounded management process.

What distinguishes the current work from the HITL frameworks surveyed
here is the role of human oversight in the process. In most HITL
systems, humans review AI output after it is generated. In the MVC
Test Orchestrator, the SRS, a human-authored document, defines the
architectural boundaries before any generation begins, so developer
intent is built into the system from the start rather than checked
after the fact.

Taken together, the literature reviewed across
Sections~\ref{sec:related_static} through~\ref{sec:related_hitl}
reveals a consistent and critical gap in the current state of the art.
Traditional static analysis tools operate at too low an abstraction
level to detect MVC-layer violations, while scenario-based evaluation
frameworks remain heavily dependent on manual configuration. LLM-based
code generation, despite its productivity gains, has been shown to
systematically undermine macro-level architectural constraints in the
absence of structured oversight. Existing traceability methods lack
the automation and dynamic integration needed to support modern
development workflows, and current automated repair mechanisms address
localized defects without accounting for broader architectural
coherence. The MVC Test Orchestrator is proposed to address this
combination of limitations. By grounding code generation in
SRS-derived requirements through RAG-based semantic retrieval,
enforcing MVC architectural boundaries through AST-based static
analysis, and embedding human oversight at each decision point through
a HITL design, the proposed system offers an integrated response to
challenges that existing approaches address only in isolation. Rather
than optimizing a single dimension of the development pipeline, the
aim is to establish a traceable, auditable, and architecturally
coherent workflow from requirements to implementation.